\begin{abstract}%
    Lahars represent one of the primary hazards associated with snow- and ice-capped volcanoes due to their ability to rapidly mobilize mixtures of water, sediment, and volcanic material through drainage networks. 
    Villarrica volcano, one of Chile's most active volcanic systems, serves as a relevant case study due to the generation of lahars during its 2015 eruptive event, which were triggered by the interaction between the 
    eruptive process and the snow and ice cover.

    This research proposes the development of a computational lahar model using Physics-Informed Neural Networks (PINNs), using the 2015 Villarrica volcano eruption as a case study. The physical model will be 
    formulated as a coupled system of partial differential equations based on the shallow water (shallow flow) approximation, incorporating mass and momentum conservation, solid concentration transport, 
    and pore pressure evolution. The PINN will approximate the spatio-temporal fields of flow depth, velocity, solid concentration, and pore pressure by embedding the governing equations directly into its loss function.

    The methodology will account for physically consistent initial and boundary conditions, topographic representation, and specific training stability strategies—including hard enforcement of initial conditions, 
    variable normalization, and convergence control to prevent trivial solutions. Subsequently, the model will be applied to the 2015 Villarrica event using available parameters and observational data. 
    The study aims to evaluate the capability of PINNs to reproduce the spatio-temporal dynamics of lahars and establish their potential as a complementary tool for the physical modeling of this volcanic hazard.
\end{abstract}

\textbf{Keywords:Physics-Informed Neural Networks (PINNs), Partial Differential Equations, Villarrica Volcano, Lahars, Volcanic Hazards} \keywords\par%